% Exercise ID: MAT_A8MODELO_1SX_NFX_011
% Module: Módulo A8 - Modelos Discretos | Concept: Sistemas Numéricos | Type: Números Figurados
% Difficulty: 3/5 (Médio) | Format: standard
% Tags: sistemas_numericos, binario, hexadecimal, sequencias, decimal, numeros_triangulares, representacao_visual, conversao_bases, regra_sucessao
% Author: Professor Diogo | Date: 2025-11-26
% Status: active

\exercicio{Considere a sucessão dos números triangulares representados visualmente abaixo (os valores $T_1$, $T_2$, $T_3$, $T_4$ estão indicados nas figuras):

\begin{center}
\begin{tikzpicture}[scale=0.5]
% T1 = 1
\foreach \x in {0} {
    \fill (\x,0) circle (0.15);
}
% Valor $T_1 = 1$ representado na figura

% T2 = 3
\begin{scope}[xshift=3cm]
\foreach \x in {0,1} {
    \fill (\x,0) circle (0.15);
}
\fill (0.5,0.5) circle (0.15);
% Valor $T_2 = 3$ representado na figura
\end{scope}

% T3 = 6
\begin{scope}[xshift=6cm]
\foreach \x in {0,1,2} {
    \fill (\x,0) circle (0.15);
}
\foreach \x in {0.5,1.5} {
    \fill (\x,0.5) circle (0.15);
}
\fill (1,1) circle (0.15);
% Valor $T_3 = 6$ representado na figura
\end{scope}

% T4 = 10
\begin{scope}[xshift=10cm]
\foreach \x in {0,1,2,3} {
    \fill (\x,0) circle (0.15);
}
\foreach \x in {0.5,1.5,2.5} {
    \fill (\x,0.5) circle (0.15);
}
\foreach \x in {1,2} {
    \fill (\x,1) circle (0.15);
}
\fill (1.5,1.5) circle (0.15);
% Valor $T_4 = 10$ representado na figura
\end{scope}
\end{tikzpicture}
\end{center}

\vspace{0.5cm}

\subexercicio{Represente os quatro primeiros termos da sucessao ($T_1, T_2, T_3, T_4$) em sistema binario.}

\subexercicio{Escreva o 5º termo da sucessao ($T_5$) em notacao decimal e hexadecimal.}

\subexercicio{Determine o 10º termo da sucessao ($T_{10}$) e represente-o em sistema binario.}

\subexercicio{Escreva a regra geral para gerar o proximo termo da sucessao dos numeros triangulares.}}
